\documentclass[11pt]{article}

\usepackage[margin=1in]{geometry}
\usepackage{amsmath,amssymb}
\usepackage{booktabs}
\usepackage{enumitem}
\usepackage[hidelinks]{hyperref}

\setlist[itemize]{leftmargin=1.4em}
\setlist[enumerate]{leftmargin=1.6em}

\title{SYSC 5108 Exam Study\\Assignment 4, Question 3}
\author{Jesse Levine}
\date{}

\begin{document}
\maketitle

\section*{Question}

Let the transition probability matrix of a two-state Markov chain be
\[
    P =
    \begin{bmatrix}
        p & 1-p \\
        1-p & p
    \end{bmatrix}.
\]

\begin{enumerate}[label=(\alph*)]
    \item Suppose at time \(0\), the distribution of states is
    \[
        Q_0 = \begin{bmatrix}0.2\\0.8\end{bmatrix}.
    \]
    Find the marginal distribution \(Q_1\) at time \(1\).
    \item Show that
    \[
        P^{(n)}
        =
        \begin{bmatrix}
            \frac12+\frac12(2p-1)^n & \frac12-\frac12(2p-1)^n \\
            \frac12-\frac12(2p-1)^n & \frac12+\frac12(2p-1)^n
        \end{bmatrix}.
    \]
    \item Find the limiting probabilities.
    \item Find
    \[
        \lim_{n\to\infty} P^{(n)}.
    \]
\end{enumerate}

\section*{Course and Textbook Cross-References}

\begin{itemize}
    \item \textbf{Chapter 6 slides, Marginal Distribution:} the state distribution
    evolves by multiplication with the transition matrix.
    \item \textbf{Chapter 6 slides, Chapman-Kolmogorov Equations:} \(n\)-step
    transition probabilities are entries of \(P^{(n)}\).
    \item \textbf{Chapter 6 slides, Continuing Cases:} limiting probabilities
    satisfy the stationary condition.
    \item \textbf{Goodfellow, Bengio, and Courville, \emph{Deep Learning},
    sequence-modeling context:} recurrent probabilistic systems evolve state
    distributions over time through repeated transition operations.
\end{itemize}

\section*{Convention}

The Chapter 6 slides write state distributions as row vectors:
\[
    \pi^T(n+1)=\pi^T(n)P.
\]
This assignment uses a \emph{column-vector} distribution \(Q_n\), so here we use
\[
    Q_{n+1}=PQ_n.
\]
The math is the same; only the orientation changes.

\section*{Part (a): Find \(Q_1\)}

Given
\[
    Q_0=
    \begin{bmatrix}
        0.2\\
        0.8
    \end{bmatrix},
\]
the distribution at time \(1\) is
\[
    Q_1 = P Q_0.
\]
Compute:
\[
\begin{aligned}
    Q_1
    &=
    \begin{bmatrix}
        p & 1-p \\
        1-p & p
    \end{bmatrix}
    \begin{bmatrix}
        0.2\\
        0.8
    \end{bmatrix} \\
    &=
    \begin{bmatrix}
        0.2p+0.8(1-p) \\
        0.2(1-p)+0.8p
    \end{bmatrix} \\
    &=
    \begin{bmatrix}
        0.8-0.6p \\
        0.2+0.6p
    \end{bmatrix}.
\end{aligned}
\]

\[
    \boxed{
    Q_1=
    \begin{bmatrix}
        0.8-0.6p\\
        0.2+0.6p
    \end{bmatrix}
    }
\]

\section*{Part (b): Show the Formula for \(P^{(n)}\)}

Define the matrices
\[
    A=
    \frac12
    \begin{bmatrix}
        1 & 1\\
        1 & 1
    \end{bmatrix},
    \qquad
    B=
    \frac12
    \begin{bmatrix}
        1 & -1\\
        -1 & 1
    \end{bmatrix}.
\]

Then
\[
    A+B = I,
    \qquad
    A-B=
    \begin{bmatrix}
        0 & 1\\
        1 & 0
    \end{bmatrix}.
\]
Using these,
\[
\begin{aligned}
    P
    &=
    \begin{bmatrix}
        p & 1-p\\
        1-p & p
    \end{bmatrix} \\
    &=
    A + (2p-1)B.
\end{aligned}
\]

Now use the algebraic properties:
\[
    A^2=A,\qquad B^2=B,\qquad AB=BA=0.
\]
Therefore,
\[
\begin{aligned}
    P^{(n)}
    &=
    \left(A+(2p-1)B\right)^n \\
    &=
    A + (2p-1)^n B,
\end{aligned}
\]
because any mixed product containing both \(A\) and \(B\) vanishes.

Substitute \(A\) and \(B\):
\[
\begin{aligned}
    P^{(n)}
    &=
    \frac12
    \begin{bmatrix}
        1 & 1\\
        1 & 1
    \end{bmatrix}
    +
    \frac{(2p-1)^n}{2}
    \begin{bmatrix}
        1 & -1\\
        -1 & 1
    \end{bmatrix} \\
    &=
    \begin{bmatrix}
        \frac12+\frac12(2p-1)^n & \frac12-\frac12(2p-1)^n \\
        \frac12-\frac12(2p-1)^n & \frac12+\frac12(2p-1)^n
    \end{bmatrix}.
\end{aligned}
\]

\[
    \boxed{
    P^{(n)}
    =
    \begin{bmatrix}
        \frac12+\frac12(2p-1)^n & \frac12-\frac12(2p-1)^n \\
        \frac12-\frac12(2p-1)^n & \frac12+\frac12(2p-1)^n
    \end{bmatrix}
    }
\]

\section*{Part (c): Find the Limiting Probabilities}

If
\[
    0<p<1,
\]
then
\[
    |2p-1|<1.
\]
Hence
\[
    \lim_{n\to\infty}(2p-1)^n = 0.
\]
So from part (b),
\[
    \lim_{n\to\infty}P^{(n)}
    =
    \begin{bmatrix}
        \frac12 & \frac12\\
        \frac12 & \frac12
    \end{bmatrix}.
\]

Therefore, for any initial distribution \(Q_0\),
\[
    \lim_{n\to\infty} Q_n
    =
    \lim_{n\to\infty} P^{(n)}Q_0
    =
    \begin{bmatrix}
        \frac12\\
        \frac12
    \end{bmatrix}.
\]

Thus the limiting probabilities are
\[
    \boxed{
    \mu=
    \begin{bmatrix}
        \frac12\\
        \frac12
    \end{bmatrix}
    \qquad (0<p<1)
    }.
\]

\subsection*{Special cases}

\begin{itemize}
    \item If \(p=1\), then
    \[
        P=I,
    \]
    so the chain never changes state, and the limit depends on the initial
    distribution:
    \[
        Q_n=Q_0.
    \]
    \item If \(p=0\), then
    \[
        P=
        \begin{bmatrix}
            0 & 1\\
            1 & 0
        \end{bmatrix},
    \]
    which swaps the states every step. The sequence oscillates, so the limit does
    not exist.
\end{itemize}

\section*{Part (d): Find \(\lim_{n\to\infty}P^{(n)}\)}

From part (b), for \(0<p<1\),
\[
    \boxed{
    \lim_{n\to\infty}P^{(n)}
    =
    \begin{bmatrix}
        \frac12 & \frac12\\
        \frac12 & \frac12
    \end{bmatrix}
    }.
\]

For the boundary cases:
\begin{itemize}
    \item \(p=1\): \(P^{(n)}=I\), so
    \[
        \lim_{n\to\infty}P^{(n)}=I.
    \]
    \item \(p=0\): \(P^{(n)}\) alternates between \(I\) and
    \(\begin{bmatrix}0&1\\1&0\end{bmatrix}\), so the limit does not exist.
\end{itemize}

\section*{Final Answer}

\[
    \boxed{
    Q_1=
    \begin{bmatrix}
        0.8-0.6p\\
        0.2+0.6p
    \end{bmatrix}
    }
\]

\[
    \boxed{
    P^{(n)}
    =
    \begin{bmatrix}
        \frac12+\frac12(2p-1)^n & \frac12-\frac12(2p-1)^n \\
        \frac12-\frac12(2p-1)^n & \frac12+\frac12(2p-1)^n
    \end{bmatrix}
    }
\]

For \(0<p<1\),
\[
    \boxed{
    \lim_{n\to\infty}Q_n
    =
    \begin{bmatrix}
        \frac12\\
        \frac12
    \end{bmatrix},
    \qquad
    \lim_{n\to\infty}P^{(n)}
    =
    \begin{bmatrix}
        \frac12 & \frac12\\
        \frac12 & \frac12
    \end{bmatrix}
    }.
\]

\section*{Exam Notes}

\begin{itemize}
    \item This chain is symmetric, so the stationary distribution is uniform:
    \[
        \mu=\left[\frac12,\frac12\right]^T.
    \]
    \item The value \((2p-1)\) controls the speed of convergence. The closer
    \(|2p-1|\) is to \(1\), the slower the convergence.
    \item The special cases matter:
    \(p=1\) gives no movement, while \(p=0\) gives perfect alternation.
\end{itemize}

\end{document}
